% Scope: Derive motion-induced phase and inflow contrast; develop phase contrast, TOF and ASL with explicit timing, calibration and physiological assumptions.
% Status: architecture only; no chapter prose or completed problems yet.
% Prerequisites: Chapters 5, 7 through 9 and 17.
% Planned worked examples: Gradient moments, velocity encoding and ASL timing.
% Planned figures: Bipolar gradients, inflow saturation and label-control timing.
% Planned challenge problems: Design a VENC experiment; distinguish transit-time bias from low perfusion.
\chapter{Flow, Perfusion and Angiography}
\label{ch:19}

\section{Motion in gradient fields}
\label{sec:19:motion-in-gradient-fields}
\subsection{Position expansion and gradient moments}
\subsection{Constant velocity, acceleration and intravoxel dispersion}
\subsection{Flow compensation and residual sensitivity}

\section{Phase-contrast imaging}
\label{sec:19:phase-contrast-imaging}
\subsection{Bipolar encoding and velocity phase}
\subsection{VENC, phase wrapping and noise}
\subsection{Background phase and multi-directional flow}

\section{Inflow and angiographic contrast}
\label{sec:19:inflow-and-angiographic-contrast}
\subsection{Time-of-flight saturation and fresh spins}
\subsection{Slice/slab geometry and venous suppression}
\subsection{Contrast-enhanced angiography and timing}

\section{Arterial spin labeling}
\label{sec:19:arterial-spin-labeling}
\subsection{Pulsed, continuous and pseudocontinuous labeling}
\subsection{Label-control subtraction and kinetic models}
\subsection{Transit time, labeling efficiency and quantification}

\section{Applications and limitations}
\label{sec:19:applications-and-limitations}
\subsection{Vascular stenosis and complex flow}
\subsection{Cardiac gating and four-dimensional flow}
\subsection{Perfusion versus macrovascular contamination}

\section{Worked examples}
\label{sec:19:worked-examples}
% Planned calculations: Gradient moments, velocity encoding and ASL timing.
\subsection{Derivation and quantitative calculation}
\subsection{Assumptions, limiting cases and scanner interpretation}

\section{Challenge problems}
\label{sec:19:challenge-problems}
% Problem design: Design a VENC experiment; distinguish transit-time bias from low perfusion.
% Use challengeproblem and stable labels prob:19:<topic>.
% Complete solutions belong in Appendix E, section for this chapter.
% Use \begin{solution}{prob:19:<topic>} ... \end{solution} there.
\subsection{Derivation and quantitative design}
\subsection{Model criticism and integrated reasoning}
